\chapter{The Head Order Reduction System}
\label{chapter-horse}

	This chapter describes HORSE, the programming environment for THOR.
The name HORSE is taken from the strategy used to implement THOR, which
is the head order graph reduction method of Berkling \cite{berkling}.  
HORSE provides
facilities for associating names with expressions, defining structures,
loading files of expressions, and controlling the reduction of expressions.

	In the sections that follow, {\tt any commands that should be typed on the
computer will be in this typeface}, and {\sc the responses of the computer
will be in this typeface}.


\section{Running HORSE}

	HORSE currently runs on IBM PC's and SUN workstations.  To invoke the
HORSE system, simply type {\tt horse} at the system level prompt.  When
HORSE is up and running, it will print a banner message which gives the
current version number and some memory size information, and then print the
HORSE system prompt, which is:
\begin{quote}
	{\sc horse>}
\end{quote}
The {\tt horse} command can be optionally followed by the names of files to
be automatically loaded into the programming environment.  For example,
the command
\begin{quote}
	{\sc \%}{\tt horse comp.lam funs.lam}
\end{quote}
will cause the files `comp.lam' and `funs.lam' to be loaded when HORSE
is brought up.  Note that there are different filename and search path
conventions on different computers, and it may be necessary to supply
full directory names for some files.


\section{A few words about input}

	HORSE itself is case insensitive;  everything that is read in is
converted to upper-case.  This can cause problems when dealing with files
on operating systems that are case sensitive.

	The semicolon character is used to mark the beginning of a comment in
HORSE and THOR.  Any text after the semicolon but before a newline
is ignored.

	The reduction process can (sometimes) be halted by typing a
control-C, which will get you back to the HORSE prompt.  On UNIX systems,
the HORSE process can be suspended by typing the usual control-Z when
at the HORSE prompt.

	An input editor is provided to help ease the keyboard entry of
expressions; it is described in detail in Section~\ref{sec-edit}.



\section{Defining Forms}
\label{sec-defining}

	There are three special defining forms in HORSE; two of them allow you
to associate an expression with a symbol, and the third assists in
defining structured data objects.

\subsection{{\tt Define}}

	The {\tt define} special form allows you to associate an expression
with a symbol. It has the syntax
\begin{quote}
	{\tt (define {\sl symbol  expr})}
\end{quote}
When a {\tt define} form is seen by horse, it associates {\sl expr} 
with {\sl symbol}, so that {\sl symbol} now reduces to
{\sl expr}.  Any previous association {\sl symbol} may have had is
forgotten.  For example,
\begin{quote}
	{\sc horse>}{\tt (define add2 (lambda (x) (+ x 2)))}
\end{quote}
defines an expression, {\tt add2}, which adds two to its argument.

	It is possible to redefine the built-in primitives of HORSE,
although this is not recommended.
When you redefine a primitive, a warning message is printed on the
screen.  If you accidentally redefine a primitive, the only thing you can
do to get the original primitive back is to exit the system, correct any
files if necessary, and reload.

\subsection{{\tt Rdefine}}

	The {\tt rdefine} special form is similar to the {\tt define} special
form, except that the expression is reduced before it is associated with
the symbol.  It has the syntax
\begin{quote}
	{\tt (rdefine {\sl symbol  expr})}
\end{quote}
The {\tt rdefine} form is extremely useful as a software engineering
tool.  An extended example of THOR/HORSE's unique abilities in this
regard is given in Chapter~\ref{chapter-examples}.


\subsection{{\tt Defstruct}}
\label{sec-defstruct}
	
	The {\tt defstruct} special form is used to define new kinds of 
compound data objects, called structures.    {\tt Defstruct} has
the syntax
\begin{quote}
	{\tt(defstruct {\sl tag ${\mbox {\sl element}}^*$})}
\end{quote}
where {\sl tag} is the type of the structure, and 
${\mbox {\sl element}}^*$ is zero or more names that identify particular
elements
within the structure.  (The term `type' is used rather loosely here,
because THOR is an untyped language.)  

	When it sees this form, HORSE defines an expression for creating 
a structure instance of type {\sl tag}, which is associated with the
symbol {\tt make-{\sl tag}}.  HORSE also defines expressions for 
selecting a structure's
elements, which are associated with the symbols {\tt{\sl tag}-{\sl
element}}.  For example, the form
\begin{quote}
	{\tt (defstruct tree label branches)}
\end{quote}
defines a structure named {\tt tree}, which has two elements, {\tt
label} and {\tt branches}.  The structure creating expression is
associated with the symbol {\tt make-tree}, and a tree's elements may be
accessed with {\tt tree-label} and {\tt tree-branches}.

	

\section{Commands}

	There are only a handful of HORSE commands, each of which will be
presented in alphabetical order.    Many of these are {\em toggle
commands\/} that set the value of certain binary switches which control how
HORSE and THOR behave.  Toggle commands can take an optional argument, which must be
either {\tt on} or {\tt off}.  If a toggle command is entered without an
argument, it's current switch setting is printed.
\begin{quote}
{\bf All horse commands should be preceded
by a comma on the command line.}
\end{quote}

\subsection{{\tt EXIT}}

	The {\tt exit} command terminates execution of HORSE and returns the user
to the host system.

\subsection{{\tt FORCE}}

	The {\tt force} command toggles whether compound data objects are reduced
lazily or eagerly.


\subsection{{\tt LOAD}}

	The {\tt load} command is used to read files into the HORSE
environment.  This command takes a single argument, which is the name of
the file.  Be aware that some file systems (UNIX in particular) are case
sensitive, and HORSE is not.  If you make sure that your file names
are either all upper case letters or all lower case letters, you
won't have any trouble.

\subsection{{\tt LOG}}

	The {\tt log} command lets you make a log file which records everything
that is printed on the screen during a HORSE session.  Logging
is started with the command {\tt log {\sl filename}}, where {\sl
filename} will be the name of the log file, and is stopped with the
command {\tt log end}.  Needless to say, don't try to name your log file
{\tt end}, or you won't get very far.


\subsection{{\tt STATS}}

	The {\tt stats} command toggles the printing of execution statistics
after an expression is reduced.

\subsection{{\tt STEP}}

	The {\tt step} command toggles whether HORSE should fully reduce an
expression or wait for the user to specify how many reductions to
perform.  When the step switch is on, the user has control over how many
reductions HORSE can perform on the {\em focus expression}.  When a new
expression is entered at the HORSE prompt, it always becomes the focus
expression.  (Note that this is only true of expressions, not defining
forms or commands.)  The focus expression may be reduced by specifying the allowable
reductions at the HORSE prompt.  This is done by entering a comma
followed by an integer.  The focus expression may be completely reduced
by entering two consecutive commas.  The reduced expression now
becomes the focus expression.

	An example of how stepping works is useful.  First, HORSE will be put
into step mode:
\begin{quote}
{\samepage
	{\sc horse>}{\tt ,step on}\\
	{\sc step enabled}
}
\end{quote}
Now we can step through a computation by
specifying the maximum number of reductions that may be performed on 
an expression at one time:
\pagebreak
\begin{quote}
{\samepage
	{\sc horse>}{\tt (+ 1 (+ 2 3))}\\
	{\sc horse>}{\tt ,1}\\
	{\sc (+ 1 5)}\\
	{\sc 1 reductions this step.}\\
	{\sc 1 reductions, total.}
}
\end{quote}
When  (+ 1 (+ 2 3)) was entered at the prompt, it became the focus
expression.  On the next line, HORSE is told to reduce this expression at
most once, yielding the expression (+ 1 5), which now becomes the
focus expression.  HORSE was kind enough to tell us the number of
reductions that were performed during this step and the total number of
reductions that have been performed on this expression since it was
first entered at the prompt.  Continuing:
\begin{quote}
{\samepage
	{\sc horse>}{\tt ,1}\\
	{\sc 6}\\
	{\sc 1 reductions this step.}\\
	{\sc 2 reductions, total.}
}
\end{quote}
You can completely reduce the focus expression by using the
double comma, as in this example:
\begin{quote}
{\samepage
	{\sc horse>}{\tt ((lambda (x) (if (< x 0) (minus x) x)) -18)}\\
	{\sc horse>}{\tt ,1}\\
	{\sc (if (< -18 0) (minus -18) -18)}\\
	{\sc 1 reductions this step.}\\
	{\sc 1 reductions, total.}\\
	{\sc horse>}{\tt ,,}\\
	{\sc 18}\\
	{\sc 3 reductions this step.}\\
	{\sc 4 reductions, total.}
}
\end{quote}

	In reality, HORSE is always in step mode, but the number of
reductions it allows when the user turns off step mode is very large.
This number is not infinity, however, and it is possible to exhaust all
of these reductions and still not be in normal form.  If this is the
case, you should just enter the double comma command to continue  (you
don't need to put HORSE in step mode for it to understand the double comma).

\section{The Input Editor}
\label{sec-edit}

	HORSE is equipped with an input editor that helps you to see the
structure of your expressions, lets you make minor changes without
too many hassles, and maintains an input and output history which may be
retrieved. 

\subsection{Fence Matching}
	To help keep track of parens, brackets, and braces,
the input editor highlights matching fence pairs as they are
entered.  If a matching fence is not found, the editor beeps.
The editor will not let you enter an expression with mismatched or
unmatched fence characters. 


\subsection{Editing Commands}

	The following is a description of the editing keystroke commands and their
functions.   Most of the keystrokes are modeled
after the Emacs command that performs the same function.
Not all of the keystrokes indicated
may be available on all systems. 

\subsubsection{Cursor Movement Commands}
\begin{tabbing}
CTRL-F, Right-arrow\ \   \=move cursor right one character\\
CTRL-Right-arrow \>move cursor to beginning of word to the right\\
CTRL-B, Left-arrow \>move cursor left one character\\
CTRL-Left-arrow \>move cursor to beginning of word to the left\\
CTRL-P, Up-arrow \>move cursor one line upwards, or beginning of line\\
CTRL-N, Down-arrow \>move cursor one line downwards, or end of line\\
CTRL-A, Home \>move cursor to beginning of line\\
CTRL-E, End \>move cursor to end of line
\end{tabbing}

\subsubsection{Deletion Commands}
\begin{tabbing}
CTRL-D, Delete\ \    \=delete the character under the cursor\\
Backspace \>delete the character to the left of cursor\\
CTRL-K \>delete all characters from cursor to end of line
\end{tabbing}

\subsection{History Commands}

	The input and output histories maintain a record of recent input
and output expressions and definitions.  (Note that commands are not kept.)
These histories
are kept in circular buffers of finite length, so only the most recent
expressions are kept.  The contents of a buffer may be inserted at the
current cursor position using the following keystrokes:
\begin{quote}
\begin{tabbing}
CTRL-Y    \=yanks a string from the input history buffer\\
CTRL-R \>yanks a string from the output history buffer
\end{tabbing}
\end{quote}
The contents of a history may be cycled through by repeating the
command keystroke;  when the end of the history is reached, an empty
string is inserted at the cursor.
